%% Template for a CV
%% Author: Rob J Hyndman

\documentclass[10pt,letterpaper,]{article}
\usepackage[scaled=0.86]{DejaVuSansMono}
\usepackage[sfdefault,lf,t]{carlito}

% Change color to blue
\usepackage{color,xcolor}
\definecolor{headcolor}{HTML}{990000}

\usepackage{ifxetex,ifluatex}
\usepackage{fixltx2e} % provides \textsubscript
\ifnum 0\ifxetex 1\fi\ifluatex 1\fi=0 % if pdftex
  \usepackage[T1]{fontenc}
  \usepackage[utf8]{inputenc}
\else % if luatex or xelatex
  \ifxetex
    \usepackage{mathspec}
  \else
    \usepackage{fontspec}
  \fi
  \defaultfontfeatures{Ligatures=TeX,Scale=MatchLowercase}
\fi

\usepackage[utf8]{inputenc}
\usepackage[T1]{fontenc}

% use upquote if available, for straight quotes in verbatim environments
\IfFileExists{upquote.sty}{\usepackage{upquote}}{}
% use microtype if available
\IfFileExists{microtype.sty}{%
\usepackage[]{microtype}
\UseMicrotypeSet[protrusion]{basicmath} % disable protrusion for tt fonts
}{}
\PassOptionsToPackage{hyphens}{url} % url is loaded by hyperref
\usepackage[unicode=true,hidelinks]{hyperref}
\urlstyle{same}  % don't use monospace font for urls
\usepackage{geometry}
\geometry{left=1.75cm,right=1.75cm,top=2.2cm,bottom=2cm}

\usepackage{longtable,booktabs}
% Fix footnotes in tables (requires footnote package)
\IfFileExists{footnote.sty}{\usepackage{footnote}\makesavenoteenv{long table}}{}
\IfFileExists{parskip.sty}{%
\usepackage{parskip}
}{% else
\setlength{\parindent}{0pt}
\setlength{\parskip}{6pt plus 2pt minus 1pt}
}
\setlength{\emergencystretch}{3em}  % prevent overfull lines
\providecommand{\tightlist}{%
  \setlength{\itemsep}{0pt}\setlength{\parskip}{0pt}}
\setcounter{secnumdepth}{0}

% set default figure placement to htbp
\makeatletter
\def\fps@figure{htbp}
\makeatother



\date{May 2026}


\usepackage{paralist,ragged2e,datetime}
\usepackage[hyphens]{url}
\usepackage{fancyhdr,enumitem,pifont}
\usepackage[compact,small,sf,bf]{titlesec}

\RaggedRight
\sloppy

% Header and footer
\pagestyle{fancy}
\makeatletter
\lhead{\sf\textcolor[gray]{0.4}{Cirruculum vitae: \@name}}
\rhead{\sf\textcolor[gray]{0.4}{\thepage}}
\cfoot{}
\def\headrule{{\color[gray]{0.4}\hrule\@height\headrulewidth\@width\headwidth \vskip-\headrulewidth}}
\makeatother

% Header box
\usepackage{tabularx}

\makeatletter
\def\name#1{\def\@name{#1}}
\def\info#1{\def\@info{#1}}
\makeatother
\newcommand{\shadebox}[3][.9]{\fcolorbox[gray]{0}{#1}{\parbox{#2}{#3}}}

\usepackage{calc}
\newlength{\headerboxwidth}
\setlength{\headerboxwidth}{\textwidth}
%\addtolength{\headerboxwidth}{0.2cm}
\makeatletter
\def\maketitle{
\thispagestyle{plain}
\vspace*{-1.4cm}
\shadebox[0.9]{\headerboxwidth}{\sf\color{headcolor}\hfil
\hbox to 0.98\textwidth{\begin{tabular}{l}
\\[-0.3cm]
\LARGE\textbf{\@name}
\\[0.1cm]\large Data Scientist\\[0.6cm]
\normalsize\textbf{Cirruculum vitae}\\
\normalsize May 2026
\end{tabular}
\hfill\hbox{\fontsize{9}{12}\sf
\begin{tabular}{@{}rl@{}}
\@info
\end{tabular}}}\hfil
}
\vspace*{0.2cm}}
\makeatother

% Section headings
\titlelabel{}
\titlespacing{\section}{0pt}{1.5ex}{0.5ex}
\titleformat*{\section}{\color{headcolor}\large\sf\bfseries}
\titleformat*{\subsection}{\color{headcolor}\sf\bfseries}
\titlespacing{\subsection}{0pt}{1ex}{0.5ex}

% Miscellaneous dimensions
\setlength{\parskip}{0ex}
\setlength{\parindent}{0em}
\setlength{\headheight}{15pt}
\setlength{\tabcolsep}{0.15cm}
\clubpenalty = 10000
\widowpenalty = 10000
\setlist{itemsep=1pt}
\setdescription{labelwidth=1.2cm,leftmargin=1.5cm,labelindent=1.5cm,font=\rm}

% Make nicer bullets
\renewcommand{\labelitemi}{\ding{228}}

\usepackage{booktabs,fontawesome5}
%\usepackage[t1,scale=0.86]{sourcecodepro}

\name{Abhijit Dasgupta}
\def\imagetop#1{\vtop{\null\hbox{#1}}}
\info{%
\raisebox{-0.05cm}{\imagetop{\faIcon{map-marker-alt}}} &  \imagetop{\begin{tabular}{@{}l@{}}12410
Milestone Manor Lane, Germantown MD 20876\end{tabular}}\\ %
% %
%
\faIcon{phone-alt} & +1 (240) 813-8458\\%
\faIcon{envelope} & \href{mailto:adasgupta@araastat.com}{\nolinkurl{adasgupta@araastat.com}}\\%
\faIcon{twitter} & \href{https://twitter.com/webbedfeet}{@webbedfeet}\\%
\faIcon{github} & \href{https://github.com/webbedfeet}{webbedfeet}\\%
\faIcon{linkedin} & \href{https://www.linkedin.com/in/abhijit-dasgupta-70197713}{abhijit-dasgupta-70197713}\\%
%
%
%
}


%\usepackage{inconsolata}


\setlength\LTleft{0pt}
\setlength\LTright{0pt}

% Pandoc CSL macros
% definitions for citeproc citations
\NewDocumentCommand\citeproctext{}{}
\NewDocumentCommand\citeproc{mm}{%
  \begingroup\def\citeproctext{#2}\cite{#1}\endgroup}
\makeatletter
 % allow citations to break across lines
 \let\@cite@ofmt\@firstofone
 % avoid brackets around text for \cite:
 \def\@biblabel#1{}
 \def\@cite#1#2{{#1\if@tempswa , #2\fi}}
\makeatother
\newlength{\cslhangindent}
\setlength{\cslhangindent}{1.5em}
\newlength{\csllabelwidth}
\setlength{\csllabelwidth}{3em}
\newenvironment{CSLReferences}[2] % #1 hanging-indent, #2 entry-spacing
 {\begin{list}{}{%
  \setlength{\itemindent}{0pt}
  \setlength{\leftmargin}{0pt}
  \setlength{\parsep}{0pt}
  % turn on hanging indent if param 1 is 1
  \ifodd #1
   \setlength{\leftmargin}{\cslhangindent}
   \setlength{\itemindent}{-1\cslhangindent}
  \fi
  % set entry spacing
  \setlength{\itemsep}{#2\baselineskip}}}
 {\end{list}}

\usepackage{calc}
\newcommand{\CSLBlock}[1]{\hfill\break\parbox[t]{\linewidth}{\strut\ignorespaces#1\strut}}
\newcommand{\CSLLeftMargin}[1]{\parbox[t]{\csllabelwidth}{\strut#1\strut}}
\newcommand{\CSLRightInline}[1]{\parbox[t]{\linewidth - \csllabelwidth}{\strut#1\strut}}
\newcommand{\CSLIndent}[1]{\hspace{\cslhangindent}#1}


\def\endfirstpage{\newpage}

\begin{document}
\maketitle


\section{Summary}\label{summary}

I am an expert statistician and data scientist interested in
multidisciplinary collaborations to understand and solve problems using
modern data analytic tools and techniques. I work at the interface of
statistics and machine learning, and have extensive experience (SME) in
statistical methods, study design, machine learning, Bayesian analyses
and data visualization. I am an expert in using R and Python for data
analytics and data science applied to a wide variety of disciplines,
with over 20 years experience. I am an excellent communicator and
teacher who can explain technical details well to a non-technical
audience, while maintaining rigor in the analytic techniques used. I am
also interested in the mentoring, education and training of the next
generation of data scientists.

\section{Current Positions}\label{current-positions}

\subsection{Primary}\label{primary}

\begin{longtable}{@{\extracolsep{\fill}}ll}
2025-- & \parbox[t]{0.85\textwidth}{%
\textbf{Data Science Director}\hfill{\footnotesize AstraZeneca}\newline
  Gaithersburg, MD\par%
  \vspace{0.1cm}\begin{minipage}{0.7\textwidth}%
\begin{itemize}%
\item Oncology biomarker development using statistics and machine learning%
\item Survival analysis, Bayesian modeling%
\end{itemize}%
\end{minipage}%
\vspace{\parsep}}\\
2020-- & \parbox[t]{0.85\textwidth}{%
\textbf{Data Science Associate Director}\hfill{\footnotesize AstraZeneca}\newline
  Gaithersburg, MD\par%
  \vspace{0.1cm}\begin{minipage}{0.7\textwidth}%
\begin{itemize}%
\item Biostatistics, machine learning and AI collaboration across therapeutic areas in oncology%
\item Survival analysis, machine learning, biomarker discovery and development%
\end{itemize}%
\end{minipage}%
\vspace{\parsep}}\\
2020-- & \parbox[t]{0.85\textwidth}{%
\textbf{Adjunct Professor}\hfill{\footnotesize Georgetown University}\newline
  Washington DC\par%
  \vspace{0.1cm}\begin{minipage}{0.7\textwidth}%
\begin{itemize}%
\item Data visualization (DSAN 5200)%
\item Big Data and Cloud Computing (DSAN 6000)%
\item Biological and Biomedical Data Science (DSAN 6150)%
\end{itemize}%
\end{minipage}%
\vspace{\parsep}}\\
\end{longtable}

\subsection{Other}\label{other}

\begin{longtable}{@{\extracolsep{\fill}}ll}
2009-- & \parbox[t]{0.85\textwidth}{%
\textbf{Director and Chief Data Scientist}\hfill{\footnotesize ARAASTAT}\newline
  Germantown, MD\par%
  \vspace{0.1cm}\begin{minipage}{0.7\textwidth}%
\begin{itemize}%
\item Statistics and Data Science consultant%
\item Statistical and predictive analytics, data visualization%
\item R and Python training for Data Science and Machine Learning%
\end{itemize}%
\end{minipage}%
\vspace{\parsep}}\\
2016-- & \parbox[t]{0.85\textwidth}{%
\textbf{Affiliate Faculty}\hfill{\footnotesize Department of Systems Biology, George Mason University}\newline
  Fairfax, VA\par%
  \empty%
\vspace{\parsep}}\\
2020-- & \parbox[t]{0.85\textwidth}{%
\textbf{Adjunct Faculty}\hfill{\footnotesize Data Science and Analytics, Georgetown University}\newline
  Washington, DC\par%
  \empty%
\vspace{\parsep}}\\
2026-- & \parbox[t]{0.85\textwidth}{%
\textbf{Adjunct Faculty}\hfill{\footnotesize Data Science, George Washington University}\newline
  Washington, DC\par%
  \empty%
\vspace{\parsep}}\\
\end{longtable}

\section{Recent Teaching}\label{recent-teaching}

I am a \href{https://www.carpentries.org}{Carpentries}-certified
instructor with extensive experience teaching data science using R and
Python to a wide variety of corporate, government and academic
audiences. My classes have included programming neophytes as well as
experienced analysts looking to enhance their skills. Workshop and
training clients include Freddie Mac, the State Department and Brookings
Institute. I regularly lead highly-rated R-based training at AstraZeneca
and help organize internal data science forums and conferences.

\begin{tabular}{ll}
  2026 -- present & \textbf{Applied Big Data Analytics}, George Washington University \\ 
  2024 -- present & \textbf{Biological and Biomedical Data Science}, Georgetown University \\ 
  2024 -- present & \textbf{Data Visualization using R}, AstraZeneca \\ 
  2021 & \href{https://www.araastat.com/BIOF085}{\textbf{Data Science using Python}}, FAES, NIH \\ 
   & \emph{A 3-day online workshop introducing data science} \\ 
  2021 & \href{https://www.araastat.com/BIOF440}{\textbf{Data Visualization using Python}}, FAES, NIH \\ 
  2020-2023 & \textbf{Practical R}, AstraZeneca \\ 
  2020 -- 2025 & \textbf{Advanced Data Visualization}, Georgetown University \\ 
  2020 -- 2024 & \textbf{Big Data Analytics and Cloud Computing}, Georgetown University \\ 
  2019 -- 2020 & \href{https://www.araastat.com/BIOF439}{\textbf{Data Visualization using R}}, FAES, NIH \\ 
  2019 & \textbf{Programming with R}, Foreign Service Institute, Department of State \\ 
   & \emph{3-day workshops introducing R in a Excel-dominant environment} \\ 
  2018 -- 2022 & \textbf{Carpentries Workshops}, NIH and other locations \\ 
   & \emph{both R- and Python-based} \\ 
  2018 & \textbf{SQL for Programmers}, Brookings Institute \\ 
  2017 & \textbf{Best Practices in Data Analysis}, Joint Statistical Meetings, Baltimore \\ 
   & \emph{This is a continuing education course. I received all A's on my evaluations} \\ 
  2017 & \textbf{Machine Learning using Python}, Freddie Mac \\ 
   & \emph{This is a 32-hour corporate training program covering the background, theory, usage } \\ 
   & \emph{and implementation of several machine learning methods} \\ 
  2016 -- 2020 & \href{https://www.araastat.com/BIOF339}{\textbf{Practical R}}, FAES, NIH \\ 
   & \emph{Introductory R course} \\ 
  \end{tabular}

\section{Education and Training}\label{education-and-training}

\begin{longtable}{@{\extracolsep{\fill}}ll}
1993 & \parbox[t]{0.85\textwidth}{%
\textbf{B.Stat.(Hons.)}\hfill{\footnotesize Indian Statistical Institute}\newline
  Kolkata, India\par%
  \empty%
\vspace{\parsep}}\\
1995 & \parbox[t]{0.85\textwidth}{%
\textbf{M.S. in Statistics}\hfill{\footnotesize University of Washington}\newline
  Seattle, WA\par%
  \empty%
\vspace{\parsep}}\\
2001 & \parbox[t]{0.85\textwidth}{%
\textbf{Ph.D. in Biostatistics}\hfill{\footnotesize University of Washington}\newline
  Seattle, WA\par%
  \empty%
\vspace{\parsep}}\\
2001-2006 & \parbox[t]{0.85\textwidth}{%
\textbf{Post-doctoral training}\hfill{\footnotesize National Cancer Institute}\newline
  Bethesda, MD\par%
  \vspace{0.1cm}\begin{minipage}{0.7\textwidth}%
\begin{itemize}%
\item Genetic epidemiology, statistical genetics, bioinformatics%
\end{itemize}%
\end{minipage}%
\vspace{\parsep}}\\
\end{longtable}

\section{Work History}\label{work-history}

\begin{longtable}{@{\extracolsep{\fill}}ll}
2018-2021 & \parbox[t]{0.85\textwidth}{%
\textbf{Co-chairman}\hfill{\footnotesize Department of Bioinformatics and Data Science, FAES Graduate School at NIH}\newline
  Bethesda, MD\par%
  \vspace{0.1cm}\begin{minipage}{0.7\textwidth}%
\begin{itemize}%
\item Oversee the data science and bioinformatics cirruculum%
\item Teach  high-demand classes on R (Practical R, and Data Visualization using R) and Python (Data Visualization using Python)%
\end{itemize}%
\end{minipage}%
\vspace{\parsep}}\\
2011-2020 & \parbox[t]{0.85\textwidth}{%
\textbf{Chief Data Scientist and Co-Founder}\hfill{\footnotesize Zansors}\newline
  Arlington, VA\par%
  \vspace{0.1cm}\begin{minipage}{0.7\textwidth}%
\begin{itemize}%
\item Lead technical efforts in R\&D, product development%
\item Algorithms, signal processing, analytics, data visualization%
\item Work with engineers, app developers, UI/UX experts, business leaders%
\end{itemize}%
\end{minipage}%
\vspace{\parsep}}\\
2010-2020 & \parbox[t]{0.85\textwidth}{%
\textbf{Biostatistician (contractor with Kelly Government Services)}\hfill{\footnotesize National Institute for Arthritis, Musculoskeletal and Skin Diseases, NIH}\newline
  Bethesda, MD\par%
  \vspace{0.1cm}\begin{minipage}{0.7\textwidth}%
\begin{itemize}%
\item Statistical and Machine Learning analyses with applications in Rheumatology%
\item Computational statistics, Bayesian methods, machine learning, simulation, survival analysis, regression, visualization%
\item Clinical trials, large administrative databases (Medicare, NIS, NHANES)%
\end{itemize}%
\end{minipage}%
\vspace{\parsep}}\\
2009-2012 & \parbox[t]{0.85\textwidth}{%
\textbf{Director, Machine Learning}\hfill{\footnotesize Tranzxn, LLC}\newline
  Gaithersburg, MD\par%
  \vspace{0.1cm}\begin{minipage}{0.7\textwidth}%
\begin{itemize}%
\item Developing models for predicting hospital acquired conditions%
\end{itemize}%
\end{minipage}%
\vspace{\parsep}}\\
2006-2009 & \parbox[t]{0.85\textwidth}{%
\textbf{Assistant Professor}\hfill{\footnotesize Division of Biostatistics, Thomas Jefferson University}\newline
  Philadelphia, PA\par%
  \vspace{0.1cm}\begin{minipage}{0.7\textwidth}%
\begin{itemize}%
\item Biostatistics, Bioinformatics, Oncology%
\item Supported the Kimmel Cancer Center%
\end{itemize}%
\end{minipage}%
\vspace{\parsep}}\\
\end{longtable}

\section{Professional Service}\label{professional-service}

\begin{enumerate}
\def\labelenumi{\arabic{enumi}.}
\tightlist
\item
  Member, Data Safety Monitoring Boards, National Institute of Mental
  Heath, NIH
\item
  Member, NIH SBIR review committees related to neurobiology
\item
  Reviewer for various academic journals.
\end{enumerate}

\section{Professional Memberships}\label{professional-memberships}

\begin{enumerate}
\def\labelenumi{\arabic{enumi}.}
\tightlist
\item
  American Statistical Association
\end{enumerate}

\section{Community Leadership}\label{community-leadership}

\begin{tabular}{ll}
  2018-2019 & \textbf{Board of Directors}, \href{http://www.sanskriti-dc.org}{Sanskriti, Inc.} \\ 
   & \emph{A socio-cultural organization supporting the Bengali community} \\ 
   & \emph{in Montgomery County and the greater DC area} \\ 
  2013-2016 & \textbf{Board of Directors}, \href{https://www.datacommunitydc.org}{Data Community DC} \\ 
  2011-2020 & \textbf{Co-organizer}, Statistical Programming DC \\ 
   & \emph{Originally the DC R Users Group, this meetup brings together} \\ 
   & \emph{the DC data science community to share and discuss data science implementations in R and Python} \\ 
  \end{tabular}

\section{Publications}\label{publications}

\subsection{Ph.D.~Thesis}\label{ph.d.-thesis}

\phantomsection\label{refs-62a838691798a1c4ca76e8035c2a30c5}
\begin{CSLReferences}{0}{0}
\bibitem[\citeproctext]{ref-Dasgupta2001}
\CSLLeftMargin{1. }%
\CSLRightInline{Dasgupta, A. (2001). \emph{Parametric Identifiability
and related problems} {[}PhD thesis{]}. University of Washington.}

\end{CSLReferences}

\subsection{Books}\label{books}

\phantomsection\label{refs-9b1a6c5914c6fcea964e6672e8920d4b}
\begin{CSLReferences}{0}{0}
\bibitem[\citeproctext]{ref-Tattar2017}
\CSLLeftMargin{1. }%
\CSLRightInline{Tattar, P., Ojeda, T., Murphy, S. P., Bengfort, B., \&
Dasgupta, A. (2017). \emph{Practical data science cookbook, second
edition}. Packt Publishing.}

\bibitem[\citeproctext]{ref-Ojeda2014}
\CSLLeftMargin{2. }%
\CSLRightInline{Ojeda, T., Murphy, S. P., Dasgupta, A., \& Bengfort, B.
(2014). \emph{Practical data science cookbook}. Packt Publishing.}

\end{CSLReferences}

\subsection{Articles}\label{articles}

I have 72 peer-reviewed publications and abstracts. A full listing and
citation statistics are available at my
\href{https://scholar.google.com/citations?user=EgaGUCwAAAAJ&hl=en}{Google
Scholar page}\\
\vspace{0.5cm}

\section{Publications}\label{publications-1}

\subsection{Ph.D.~Thesis}\label{ph.d.-thesis-1}

\subsection{Books}\label{books-1}

\subsection{Articles}\label{articles-1}

\subsection{Abstracts}\label{abstracts}

\subsection{Manuscripts in progress}\label{manuscripts-in-progress}

--\textgreater{}

\end{document}
